\documentclass[10pt]{article}
\usepackage[a4paper,margin=15mm]{geometry}
\usepackage{xcolor}
\usepackage{modernnewspaper}
\definecolor{newsblue}{HTML}{164A70}
\definecolor{newscoral}{HTML}{D2644A}
\SetPaperName{The Common Thread}
\SetPaperSlogan{Clear news for curious neighbours}
\SetPaperLocation{Your town}
\SetPaperWebsite{example.org/thread}
\SetPaperVolume{4}
\SetPaperIssue{12}
\SetPaperDate{August 2026}
\renewcommand{\headlinestyle}{\large\sffamily\bfseries\color{newsblue}}
\renewcommand{\bylinestyle}{\small\sffamily\color{newscoral}}
\renewcommand{\datelinestyle}{\scriptsize\sffamily}
\begin{document}
\MakePaperHeader
\BeginNewsColumns{3}
\headline{A better route to the market}
\byline{The Field Notes Desk}
\dateline{August 2026}
\DropCap{T}he most useful local updates are short enough to read at the
corner shop and specific enough to change what happens next. This issue
follows one walking route, three observations, and a small experiment with
shade markers.

\headline{Three things we noticed}
\begin{itemize}
  \item Shade and a seat work best within a few steps of each other.
  \item A clear route marker prompts more questions than a dense map.
  \item The busiest crossing is also the place where the queue moves slowest.
\end{itemize}

\headline{What comes next}
Repeat the observation after rain. Bring one unanswered question and one idea
small enough to try before the following issue.

\headline{Noticeboard}
\textbf{Open walk}\par
Wednesday, 18:00 --- meet by the library steps for a 30-minute route review.

\medskip
Send a short observation, a date, and a place to
\href{mailto:notes@example.org}{notes@example.org}.
\EndNewsColumns
\end{document}
